\section{Independent NSFG outcome and denominator audit}
The original female respondent file contained 5,586 records; the strict interview-age 20--49 domain contained 4,855. Including one interview-age-50 respondent eligible at screening gave 4,856. Valid diagnosis groups were analysed separately, excluding unknown responses. Broad fertility problems were FECUND in 3--5 or INFERT=2; surgical sterility was excluded from impaired fecundity. Current union was RMARITAL in 1--2. The narrow outcome was INFERT=2, retaining valid INFERT categories 1--3 in the denominator. Narrow diagnosed numerator/domain counts were 15/173 for endometriosis and 29/252 for PCOS.

For ratio estimate $p_j=\sum_iw_in_{ij}/\sum_iw_id_{ij}$ the linearized respondent vector is $u_{ij}=w_i(n_{ij}-p_jd_{ij})/\sum_iw_id_{ij}$. Sum vectors by PSU within stratum and compute
\[\widehat V=\sum_h\frac{m_h}{m_h-1}\sum_b(\boldsymbol U_{hb}-\overline{\boldsymbol U}_h)(\boldsymbol U_{hb}-\overline{\boldsymbol U}_h)^T.\]
All 80 PSUs in 20 strata were retained. The 36-by-36 matrix preserves dependence across outcome definitions. R survey and Python reconstructions agreed within approximately $10^{-15}$ for estimates/standard errors and $10^{-18}$ for covariance. These are computational checks, not GBD validation. Pointwise $t_{60}$ intervals do not establish familywise significance or causal effects.

PCOS prevalence-difference shifts were 5.43 percentage points (95\% CI 0.33--10.53) for restricting the broad-outcome population to unions, and 17.15 (10.24--24.06) for changing broad to 12-month outcome within unions. The total shift was 22.58 (16.07--29.08), using joint covariance. Endometriosis shifts were 8.28 (3.01--13.56), 7.34 ($-1.44$ to 16.11) and 15.62 (7.32--23.92). All are exploratory and unadjusted.

Ever fertility services were HLPPRG=1 or HLPMC=1, with definite negatives requiring both responses=5. Complete-case estimates were 30.84/12.72\% for endometriosis diagnosed/not diagnosed and 34.49/11.80\% for PCOS. They did not all equal the CDC rounded percentages 30.7/12.7 and 34.5/11.7. The report does not provide the full missing-response implementation; the discrepancy is not hidden or forced to agree. Six core published percentages and two diagnosis standard errors did agree to published precision.

\section{Annual-chained decomposition sensitivity}
The endpoint product decomposition can be applied to each adjacent year and summed. The PCOS reference-share contribution becomes +87.11 per 100,000 (57.5\%), compared with endpoint +88.65 (58.5\%). Both sum with their rate components to +151.54. The Chinese-rate component remains +271.51; the chained reference-rate component is $-207.08$. The decomposition describes time-path allocation, not a new method or GBD uncertainty.

\section{GBD joint uncertainty remains unavailable}
The IHME request received an automatic acknowledgement, not data or approval. Original marginal intervals do not identify cross-cell dependence. NSFG covariance cannot replace jointly indexed GBD draws. GBD derived results remain point estimates. The survey is not a proven held-out source relative to GBD and does not validate the Chinese endometriosis terminal-year change.
